\documentclass[10pt,a4paper]{amsart}
\usepackage[margin=19mm]{geometry}
\usepackage{amsmath,amssymb,mathtools}
\usepackage[scr=rsfs,cal=euler]{mathalfa}
\usepackage{xfrac}
\usepackage[unicode,hidelinks,bookmarks=false]{hyperref}
\newcommand{\ie}{i.e.\ }
\newcommand{\cf}{cf.\ }
\newcommand{\resp}{respectively\ }
\newcommand{\Subsection}{\subsection}
\providecommand{\nobreakdash}{\nobreak\hbox{-}}
\setlength{\emergencystretch}{2em}
\multlinegap0pt
\usepackage{amssymb}
\usepackage{mathtools}
\usepackage[shortlabels]{enumitem}
\usepackage{xcolor}
\newtheorem{proposition}{Proposition}
\newcommand{\R}{\mathbb R}
\newcommand{\abs}[1]{\lvert #1\rvert}
\newcommand{\dd}{\,\mathrm d}
\newcommand{\dsec}{\delta_{\mathrm{sec}}}
\newcommand{\eps}{\varepsilon}
\title{Curvature Sign Rigidity and Sharp Pointwise Pinching Thresholds}
\author{Minbo Gao, Yuhang Liu, Genyuan Zhang}
\date{Source: arXiv:2609.07252v1 (CC BY 4.0)}
\begin{document}
\maketitle

\noindent\emph{Source and licence.} Curvature Sign Rigidity and Sharp Pointwise Pinching Thresholds. Version arXiv:2609.07252v1 (CC BY 4.0), distributed by arXiv under the Creative Commons Attribution 4.0 International licence, http://creativecommons.org/licenses/by/4.0/, as recorded in the arXiv licence metadata for this exact version and shown on the abstract page https://arxiv.org/abs/2609.07252v1. Full licence notice: \texttt{licenses/arxiv-2609.07252-CC-BY-4.0.txt}.\par\smallskip

\noindent\textit{Editorial scope.} Counted unit: the article's proposition prop:conformal-signchange (source0.tex lines 757--767). The article's complete original proof of it is delivered with it (source0.tex lines 769--849, one \texttt{proof} environment of 81 lines). No mathematics is written, repaired or re-derived: the statement and the proof are the article's own lines, in the article's own notation, and no retained line is substituted or re-flowed.  Retained uncounted as the source-local context the proof needs: lines 752--754.  Nothing was omitted from any retained span, so the package carries no omission record.  No retained line contains a citation, so the wrapper carries no bibliography and no bibliography entry.  No retained line contains a figure environment, an image-inclusion command or drawing code, so no image is delivered and the package declares no asset dependency.  Editorial formatting only, in addition: an AMS \texttt{amsart} preamble that restates the article's own declarations and macros used by the retained lines, namely the environments \texttt{proposition} on separate counters, together with the standard packages those lines need.  The retained environments print in this document as Proposition 1 in order of appearance, whereas the article prints the same items with the numbering of its own full document.  The complete native source of the article as distributed in the arXiv e-print for this exact version is bundled unchanged as \texttt{sources/209613/source0.tex}.
\begin{equation}\label{eq:Schouten-conf}
S=-D^2u+\dd u\otimes\dd u-\frac12\abs{Du}^2g_{\mathrm{Euc}}.
\end{equation}

\begin{proposition}[Degenerate-pinching sign change]\label{prop:conformal-signchange}
For every $n\ge3$ there is a connected smooth conformally flat metric on a slab in $\R^n$ satisfying the sectional sign trichotomy, with negative sectional curvature on one side of a flat hypersurface and positive sectional curvature on the other.  If $d$ denotes distance to the hypersurface, then
\begin{equation}\label{eq:degenerate-rates}
\dsec=e^{-1/d+o(1)}
\quad\text{on the negative side},
\qquad
\dsec=(4+o(1))d^4
\quad\text{on the positive side}.
\end{equation}
The asymptotics are uniform in the transverse variable on a fixed smaller slab.
\end{proposition}

\begin{proof}
Write $x=(t,z)\in\R\times\R^{n-1}$ and define the one-sided flat functions
\[
F(t)=\begin{cases}e^{1/t},&t<0,\\0,&t\ge0,\end{cases}
\qquad
H(t)=\begin{cases}e^{-1/t},&t>0,\\0,&t\le0.\end{cases}
\]
For small fixed $\eps>0$, set
\begin{equation}\label{eq:u-conformal}
u(t,z)=F(t)-\eps H(t)e^{\abs z^2},
\qquad
g=e^{2u}g_{\mathrm{Euc}}.
\end{equation}
We restrict to $-1/2<t<t_0$ and $\abs z<r_0$, with $t_0,r_0$ sufficiently small.

For $t<0$, $u=F(t)$ depends only on $t$.  The two types of eigenvalue sums in \eqref{eq:Schouten-conf} are
\[
-F''(t)
\quad\text{and}\quad
-(F'(t))^2.
\]
Here
\[
F'(t)=-\frac{e^{1/t}}{t^2},
\qquad
F''(t)=\frac{e^{1/t}(1+2t)}{t^4}>0
\quad(-1/2<t<0).
\]
Thus all sectional curvatures are negative.  Near $t=0$ the smaller magnitude is $(F')^2$ and the larger is $F''$, so
\[
\dsec(t,z)=\frac{(F')^2}{F''}
=\frac{e^{1/t}}{1+2t}
=e^{-1/\abs t+o(1)}.
\]

For $t>0$, put
\[
q=\eps e^{-1/t}e^{\abs z^2},
\qquad u=-q.
\]
Relative to the splitting $\R\partial_t\oplus\R^{n-1}$, the matrix of $S$ is
\[
S=\begin{pmatrix}a&b^T\\ b&C\end{pmatrix},
\]
where, with $s=\abs z^2$,
\begin{align*}
a&=\frac{q(1-2t)}{t^4}+\frac{q^2}{2t^4}-2q^2s,\\
b&=\frac{2q(1+q)}{t^2}z,\\
C&=\left(2q-\frac{q^2}{2t^4}-2q^2s\right)I
  +4q(1+q)z\otimes z.
\end{align*}
Because $q/t^4\to0$ as $t\downarrow0$, uniformly for $\abs z\le r_0$, one has
\[
\frac Cq\longrightarrow 2I+4z\otimes z,
\qquad
\frac{t^4a}{q}\longrightarrow1,
\qquad
\frac{t^2b}{q}\longrightarrow2z.
\]
The convergence is uniform for $\abs z\le r_0$, and therefore
\[
\frac{t^4}{q}\bigl(a-b^TC^{-1}b\bigr)
\longrightarrow
1-4z^T(2I+4z\otimes z)^{-1}z
=\frac1{1+2\abs z^2}>0.
\]
After decreasing $t_0$, the matrix $C$ and the Schur complement are positive.  Thus $S$ is positive definite and all sectional curvatures are positive.  Moreover,
\[
a^{-1}bb^T=4q\,z\otimes z+o(q),
\qquad
C-a^{-1}bb^T=2qI+o(q)
\]
uniformly for $\abs z\le r_0$.  Standard block perturbation therefore shows that the eigenvalues of $S$ consist of one eigenvalue
\[
\frac q{t^4}(1+o(1))
\]
and $n-1$ eigenvalues $2q(1+o(1))$, uniformly in $z$.  Hence the smallest sectional curvature is the sum of two transverse eigenvalues, $(4+o(1))q$, while the largest is $(1+o(1))q/t^4$.  This gives the second asymptotic in terms of the coordinate $t$.

All derivatives of $u$ vanish at $t=0$, so the metric is flat there and glues smoothly.  Moreover $u=o(1)$ uniformly in the transverse variable, and comparison with the Euclidean metric gives $d=\abs t(1+o(1))$ for the distance to the interface.  The two estimates above therefore give \eqref{eq:degenerate-rates} in terms of $d$.

\end{proof}

\end{document}
